\chapter{Ambiente Computacional e Reprodutibilidade} \label{ap:reprodutibilidade}

Em consonância com o compromisso de reprodutibilidade assumido nesta dissertação, este apêndice registra o ambiente computacional, as bibliotecas e as decisões de implementação que permitem a outro pesquisador reproduzir os resultados.

\section{Linguagem e ambiente}

Todo o \textit{pipeline} foi implementado na linguagem Python, versão 3.11, executado localmente em um ambiente isolado gerenciado pela ferramenta Conda. A opção pela execução local, e não em nuvem, decorreu da necessidade de controle sobre as versões das bibliotecas e da estabilidade do ambiente de inferência do modelo de linguagem.

\section{Bibliotecas principais}

A Tabela~\ref{tab:bibliotecas} relaciona as bibliotecas centrais empregadas em cada etapa do método, com a sua respectiva função.

\begin{table}[htpb]
\centering
\caption{Bibliotecas centrais do \textit{pipeline} e suas funções.}
\label{tab:bibliotecas}
\begin{tabular}{lll}
\hline
\textbf{Etapa} & \textbf{Biblioteca} & \textbf{Função} \\ \hline
Dados financeiros & yfinance & Coleta das cotações da PETR4 \\
Coleta de notícias & requests & Acesso à API REST dos portais \\
Sentimento & transformers & Inferência do FinBERT-PT-BR \\
Sentimento & torch (PyTorch) & Motor de execução da rede neural \\
Risco & arch & Estimação do modelo GARCH \\
Classificação & scikit-learn & Implementação do SVM e métricas \\
Classificação & xgboost & Implementação do XGBoost \\
Manipulação de dados & pandas, numpy & Processamento tabular e numérico \\
Testes estatísticos & scipy, statsmodels & Binomial, McNemar e Granger \\
Visualização & matplotlib & Geração das figuras \\ \hline
\end{tabular}
\vspace{0.2cm}
{\small \\ Fonte: Elaborado pelo autor.}
\end{table}

\section{Decisões de implementação relevantes}

Algumas decisões técnicas foram necessárias para a estabilidade e a correção do \textit{pipeline}, e ficam aqui registradas. A biblioteca \textit{transformers} foi fixada em uma versão compatível com o formato dos pesos do FinBERT-PT-BR. A semente dos geradores de números aleatórios foi fixada em um valor constante em todas as etapas que envolvem aleatoriedade, o que garante que execuções sucessivas produzam resultados idênticos. A divisão dos dados em treino, validação e teste é determinística, pois segue a ordem cronológica e não depende de sorteio. O alinhamento temporal das notícias utiliza o horário de Brasília como referência para o corte das dezessete horas, em conformidade com o horário de funcionamento da B3.

\section{Disponibilidade}

O código-fonte do \textit{pipeline}, incluindo os roteiros de coleta, de processamento de sentimento, de modelagem e de análise, foi versionado em repositório de controle de versão, o que assegura o registro histórico das alterações e a possibilidade de auditoria. Os dados brutos coletados são preservados em sua forma original, e as bases derivadas são geradas a partir deles por roteiros reproduzíveis, sem edição manual.
